% --- Archon physics formalization source begin ---
% archon:physics
% archon:covers PhyXMiniProblems/problem_phyx_mini_0573.lean
% archon:source-report reports/phyx_mini/problem_phyx_mini_0573.source.json

\chapter{Physics problem phyx\_mini\_0573}
\label{ch:PhyXMiniProblems_problem_phyx_mini_0573}

\paragraph{Problem source.}
\#\# Physical scenario

Figure shows an early proposal for a hydrogen bomb. The fusion fuel is deuterium, \textbackslash{}( \{\}\textasciicircum{}\{2\}\textbackslash{}text\{H\} \textbackslash{}). The high temperature and particle density needed for fusion are provided by an atomic bomb "trigger" that involves a \textbackslash{}( \{\}\textasciicircum{}\{235\}\textbackslash{}text\{U\} \textbackslash{}) or \textbackslash{}( \{\}\textasciicircum{}\{239\}\textbackslash{}text\{Pu\} \textbackslash{}) fission fuel arranged to impress an imploding, compressive shock wave on the deuterium. The fusion reaction is \textbackslash{}[ 5\textbackslash{} \{\}\textasciicircum{}\{2\}\textbackslash{}text\{H\} \textbackslash{}rightarrow 3\textbackslash{} \{\}\textasciicircum{}\{3\}\textbackslash{}text\{He\} + 4\textbackslash{} \{\}\textasciicircum{}\{4\}\textbackslash{}text\{He\} + 1\textbackslash{} \{\}\textasciicircum{}\{1\}\textbackslash{}text\{H\} + 2\textbackslash{}text\{n\}. \textbackslash{}]

\#\# Question

Calculate the rating of the fusion part of the bomb if it contains \$500 \textbackslash{}, \textbackslash{}text\{kg\}\$ of deuterium, \$30.0\textbackslash{}\%\$ of which undergoes fusion.

\#\# Answer choices

- A: A: 6.21 megatons of TNT

- B: B: 12.9 megatons of TNT

- C: C: 8.65 megatons of TNT

- D: D: 5.30 megatons of TNT

\#\# Figure caption (auxiliary; use the image as primary evidence)

The image depicts a diagram with two concentric circles, each centered around a black dot. The inner circle is a medium shade of blue, while the outer circle is a lighter blue. Lines extend from each black dot to text labels. 

The label for the outer circle reads "²³⁵U or ²³⁹Pu," indicating that it represents elements Uranium-235 or Plutonium-239. The label for the inner circle reads "²H," indicating it represents Deuterium, an isotope of hydrogen. 

Overall, the diagram appears to represent a conceptual model, possibly related to nuclear physics or chemistry, illustrating the relationship between these elements or isotopes.

\#\# Dataset metadata

Nuclear Physics; Physical Model Grounding Reasoning, Multi-Formula Reasoning

\paragraph{Recorded answer/context.}
C: C: 8.65 megatons of TNT

\paragraph{Figure/image path.}
/root/proposal\_for\_physic/science-mango/phyx\_mini\_run/phyx\_data/test\_image/573.png

\paragraph{Formalization target.}
create a compiling Lean file with sorry bodies at `PhyXMiniProblems/problem\_phyx\_mini\_0573.lean`. The Lean declarations must preserve the physical quantities, dimensions or dimensional roles, figure labels, governing-law hypotheses, and final relation expressed by this problem.
Use Mathlib/Physlib names found through LeanExplore where available. If a physics API is missing, introduce faithful local abstractions rather than scalar placeholder aliases.

\begin{theorem}[Physics formalization target]
\label{thm:physics:phyx_mini_0573:target}
\lean{PhyXMiniProblems.ProblemPhyXMini0573.problem_phyx_mini_0573}
\uses{def:physics:phyx-mini-0573:phyxminiproblems-problemphyxmini0573-atomicmassunitinkilograms, def:physics:phyx-mini-0573:phyxminiproblems-problemphyxmini0573-deuteriumfusionbombsetup, def:physics:phyx-mini-0573:phyxminiproblems-problemphyxmini0573-matchesfusionbombproblemdata, def:physics:phyx-mini-0573:phyxminiproblems-problemphyxmini0573-matchesprimarybombfigure, def:physics:phyx-mini-0573:phyxminiproblems-problemphyxmini0573-matchesreferencefusiondata, def:physics:phyx-mini-0573:phyxminiproblems-problemphyxmini0573-hasphysicalfusionparameters, def:physics:phyx-mini-0573:phyxminiproblems-problemphyxmini0573-satisfiesfusionyieldlaws, def:physics:phyx-mini-0573:phyxminiproblems-problemphyxmini0573-recordedanswerchoice, def:physics:phyx-mini-0573:phyxminiproblems-problemphyxmini0573-matchesdisplayedrating}
The assigned autoformalize agent should translate this physics problem into Lean declarations in the covered file, with theorem and lemma proof bodies written as `by sorry`.
\end{theorem}
\begin{proof}
This is an autoformalization task, not a proof task. Produce statements that can later be proved by the physics prover without weakening the physical model.
\end{proof}

% --- Archon declaration topology begin ---
\section{Declaration topology}
\begin{definition}[Mass Quantity]
  \label{def:physics:phyx-mini-0573:phyxminiproblems-problemphyxmini0573-massquantity}
  \lean{PhyXMiniProblems.ProblemPhyXMini0573.MassQuantity}
  A nonnegative, unit-independent physical mass
\end{definition}
\begin{proof}
  This is the declared model definition of Mass Quantity.
\end{proof}
\begin{definition}[Speed Quantity]
  \label{def:physics:phyx-mini-0573:phyxminiproblems-problemphyxmini0573-speedquantity}
  \lean{PhyXMiniProblems.ProblemPhyXMini0573.SpeedQuantity}
  A unit-independent physical speed
\end{definition}
\begin{proof}
  This is the declared model definition of Speed Quantity.
\end{proof}
\begin{definition}[mass In Kilograms]
  \label{def:physics:phyx-mini-0573:phyxminiproblems-problemphyxmini0573-massinkilograms}
  \lean{PhyXMiniProblems.ProblemPhyXMini0573.massInKilograms}
  \uses{def:physics:phyx-mini-0573:phyxminiproblems-problemphyxmini0573-massquantity}
  Kilogram readout of a physical mass in coherent SI units
\end{definition}
\begin{proof}
  This is the declared model definition of mass In Kilograms.
\end{proof}
\begin{definition}[energy In Joules]
  \label{def:physics:phyx-mini-0573:phyxminiproblems-problemphyxmini0573-energyinjoules}
  \lean{PhyXMiniProblems.ProblemPhyXMini0573.energyInJoules}
  Joule readout of a physical energy, grounded by Physlib's joule
\end{definition}
\begin{proof}
  This is the declared model definition of energy In Joules.
\end{proof}
\begin{definition}[speed In Meters Per Second]
  \label{def:physics:phyx-mini-0573:phyxminiproblems-problemphyxmini0573-speedinmeterspersecond}
  \lean{PhyXMiniProblems.ProblemPhyXMini0573.speedInMetersPerSecond}
  \uses{def:physics:phyx-mini-0573:phyxminiproblems-problemphyxmini0573-speedquantity}
  Metre-per-second readout of a physical speed in coherent SI units
\end{definition}
\begin{proof}
  This is the declared model definition of speed In Meters Per Second.
\end{proof}
\begin{definition}[exact Speed Of Light In Meters Per Second]
  \label{def:physics:phyx-mini-0573:phyxminiproblems-problemphyxmini0573-exactspeedoflightinmeterspersecond}
  \lean{PhyXMiniProblems.ProblemPhyXMini0573.exactSpeedOfLightInMetersPerSecond}
  Physlib's exact vacuum speed of light, read in metres per second
\end{definition}
\begin{proof}
  This is the declared model definition of exact Speed Of Light In Meters Per Second.
\end{proof}
\begin{definition}[atomic Mass Unit In Kilograms]
  \label{def:physics:phyx-mini-0573:phyxminiproblems-problemphyxmini0573-atomicmassunitinkilograms}
  \lean{PhyXMiniProblems.ProblemPhyXMini0573.atomicMassUnitInKilograms}
  The exact SI mass represented by one unified atomic mass unit
\end{definition}
\begin{proof}
  This is the declared model definition of atomic Mass Unit In Kilograms.
\end{proof}
\begin{definition}[mass In Atomic Mass Units]
  \label{def:physics:phyx-mini-0573:phyxminiproblems-problemphyxmini0573-massinatomicmassunits}
  \lean{PhyXMiniProblems.ProblemPhyXMini0573.massInAtomicMassUnits}
  \uses{def:physics:phyx-mini-0573:phyxminiproblems-problemphyxmini0573-massquantity, def:physics:phyx-mini-0573:phyxminiproblems-problemphyxmini0573-massinkilograms, def:physics:phyx-mini-0573:phyxminiproblems-problemphyxmini0573-atomicmassunitinkilograms}
  Atomic-mass-unit readout of a physical mass
\end{definition}
\begin{proof}
  This is the declared model definition of mass In Atomic Mass Units.
\end{proof}
\begin{definition}[Nuclear Species]
  \label{def:physics:phyx-mini-0573:phyxminiproblems-problemphyxmini0573-nuclearspecies}
  \lean{PhyXMiniProblems.ProblemPhyXMini0573.NuclearSpecies}
  Nuclear species named by the reaction, prose, or primary figure
\end{definition}
\begin{proof}
  This is the declared model definition of Nuclear Species.
\end{proof}
\begin{definition}[proton Number]
  \label{def:physics:phyx-mini-0573:phyxminiproblems-problemphyxmini0573-nuclearspecies-protonnumber}
  \lean{PhyXMiniProblems.ProblemPhyXMini0573.NuclearSpecies.protonNumber}
  \uses{def:physics:phyx-mini-0573:phyxminiproblems-problemphyxmini0573-nuclearspecies}
  Proton number of each named isotope or free neutron
\end{definition}
\begin{proof}
  This is the declared model definition of proton Number.
\end{proof}
\begin{definition}[mass Number]
  \label{def:physics:phyx-mini-0573:phyxminiproblems-problemphyxmini0573-nuclearspecies-massnumber}
  \lean{PhyXMiniProblems.ProblemPhyXMini0573.NuclearSpecies.massNumber}
  \uses{def:physics:phyx-mini-0573:phyxminiproblems-problemphyxmini0573-nuclearspecies}
  Nucleon (mass) number of each named isotope or free neutron
\end{definition}
\begin{proof}
  This is the declared model definition of mass Number.
\end{proof}
\begin{definition}[Nuclear Reaction]
  \label{def:physics:phyx-mini-0573:phyxminiproblems-problemphyxmini0573-nuclearreaction}
  \lean{PhyXMiniProblems.ProblemPhyXMini0573.NuclearReaction}
  \uses{def:physics:phyx-mini-0573:phyxminiproblems-problemphyxmini0573-nuclearspecies}
  Stoichiometric multiplicities on the two sides of a nuclear reaction
\end{definition}
\begin{proof}
  This is the declared model definition of Nuclear Reaction.
\end{proof}
\begin{definition}[five Deuterium Fusion Reaction]
  \label{def:physics:phyx-mini-0573:phyxminiproblems-problemphyxmini0573-fivedeuteriumfusionreaction}
  \lean{PhyXMiniProblems.ProblemPhyXMini0573.fiveDeuteriumFusionReaction}
  \uses{def:physics:phyx-mini-0573:phyxminiproblems-problemphyxmini0573-nuclearreaction}
  The five-deuteron fusion channel stated in the problem
\end{definition}
\begin{proof}
  This is the declared model definition of five Deuterium Fusion Reaction.
\end{proof}
\begin{definition}[Bomb Region]
  \label{def:physics:phyx-mini-0573:phyxminiproblems-problemphyxmini0573-bombregion}
  \lean{PhyXMiniProblems.ProblemPhyXMini0573.BombRegion}
  The two nested regions drawn in the primary image
\end{definition}
\begin{proof}
  This is the declared model definition of Bomb Region.
\end{proof}
\begin{definition}[Circle Shade]
  \label{def:physics:phyx-mini-0573:phyxminiproblems-problemphyxmini0573-circleshade}
  \lean{PhyXMiniProblems.ProblemPhyXMini0573.CircleShade}
  The two blue shades used to distinguish the nested circular regions
\end{definition}
\begin{proof}
  This is the declared model definition of Circle Shade.
\end{proof}
\begin{definition}[Trigger Mechanism]
  \label{def:physics:phyx-mini-0573:phyxminiproblems-problemphyxmini0573-triggermechanism}
  \lean{PhyXMiniProblems.ProblemPhyXMini0573.TriggerMechanism}
  The physical role played by the trigger according to the prose
\end{definition}
\begin{proof}
  This is the declared model definition of Trigger Mechanism.
\end{proof}
\begin{definition}[Early Fusion Bomb Figure]
  \label{def:physics:phyx-mini-0573:phyxminiproblems-problemphyxmini0573-earlyfusionbombfigure}
  \lean{PhyXMiniProblems.ProblemPhyXMini0573.EarlyFusionBombFigure}
  \uses{def:physics:phyx-mini-0573:phyxminiproblems-problemphyxmini0573-bombregion, def:physics:phyx-mini-0573:phyxminiproblems-problemphyxmini0573-circleshade}
  Qualitative information transcribed from the primary raster. The image gives no calibrated radius or length, so its geometry is categorical rather than a scalar measurement
\end{definition}
\begin{proof}
  This is the declared model definition of Early Fusion Bomb Figure.
\end{proof}
\begin{definition}[Deuterium Fusion Bomb Setup]
  \label{def:physics:phyx-mini-0573:phyxminiproblems-problemphyxmini0573-deuteriumfusionbombsetup}
  \lean{PhyXMiniProblems.ProblemPhyXMini0573.DeuteriumFusionBombSetup}
  \uses{def:physics:phyx-mini-0573:phyxminiproblems-problemphyxmini0573-massquantity, def:physics:phyx-mini-0573:phyxminiproblems-problemphyxmini0573-speedquantity, def:physics:phyx-mini-0573:phyxminiproblems-problemphyxmini0573-nuclearspecies, def:physics:phyx-mini-0573:phyxminiproblems-problemphyxmini0573-nuclearreaction, def:physics:phyx-mini-0573:phyxminiproblems-problemphyxmini0573-triggermechanism, def:physics:phyx-mini-0573:phyxminiproblems-problemphyxmini0573-earlyfusionbombfigure}
  All physical quantities needed by the yield calculation are independent fields. In particular, the fusion rating is not defined to be an answer choice or a numerical constant
\end{definition}
\begin{proof}
  This is the declared model definition of Deuterium Fusion Bomb Setup.
\end{proof}
\begin{definition}[Matches Fusion Bomb Problem Data]
  \label{def:physics:phyx-mini-0573:phyxminiproblems-problemphyxmini0573-matchesfusionbombproblemdata}
  \lean{PhyXMiniProblems.ProblemPhyXMini0573.MatchesFusionBombProblemData}
  \uses{def:physics:phyx-mini-0573:phyxminiproblems-problemphyxmini0573-massinkilograms, def:physics:phyx-mini-0573:phyxminiproblems-problemphyxmini0573-fivedeuteriumfusionreaction, def:physics:phyx-mini-0573:phyxminiproblems-problemphyxmini0573-deuteriumfusionbombsetup}
  Numerical and categorical data explicitly supplied by the problem text
\end{definition}
\begin{proof}
  This is the declared model definition of Matches Fusion Bomb Problem Data.
\end{proof}
\begin{definition}[Matches Primary Bomb Figure]
  \label{def:physics:phyx-mini-0573:phyxminiproblems-problemphyxmini0573-matchesprimarybombfigure}
  \lean{PhyXMiniProblems.ProblemPhyXMini0573.MatchesPrimaryBombFigure}
  \uses{def:physics:phyx-mini-0573:phyxminiproblems-problemphyxmini0573-deuteriumfusionbombsetup}
  Labels, colors, and nesting read directly from the supplied raster
\end{definition}
\begin{proof}
  This is the declared model definition of Matches Primary Bomb Figure.
\end{proof}
\begin{definition}[Matches Reference Fusion Data]
  \label{def:physics:phyx-mini-0573:phyxminiproblems-problemphyxmini0573-matchesreferencefusiondata}
  \lean{PhyXMiniProblems.ProblemPhyXMini0573.MatchesReferenceFusionData}
  \uses{def:physics:phyx-mini-0573:phyxminiproblems-problemphyxmini0573-energyinjoules, def:physics:phyx-mini-0573:phyxminiproblems-problemphyxmini0573-speedinmeterspersecond, def:physics:phyx-mini-0573:phyxminiproblems-problemphyxmini0573-exactspeedoflightinmeterspersecond, def:physics:phyx-mini-0573:phyxminiproblems-problemphyxmini0573-massinatomicmassunits, def:physics:phyx-mini-0573:phyxminiproblems-problemphyxmini0573-deuteriumfusionbombsetup}
  Reference values used by the recorded textbook calculation. The reaction mass defect 0.027 u is a nuclear-data input, not the requested bomb rating. The working value 3.00 × 10\textasciicircum{}8 m/s records the precision used in that calculation; Physlib's exact DimSpeed.speedOfLight remains available
\end{definition}
\begin{proof}
  This is the declared model definition of Matches Reference Fusion Data.
\end{proof}
\begin{definition}[Has Physical Fusion Parameters]
  \label{def:physics:phyx-mini-0573:phyxminiproblems-problemphyxmini0573-hasphysicalfusionparameters}
  \lean{PhyXMiniProblems.ProblemPhyXMini0573.HasPhysicalFusionParameters}
  \uses{def:physics:phyx-mini-0573:phyxminiproblems-problemphyxmini0573-massinkilograms, def:physics:phyx-mini-0573:phyxminiproblems-problemphyxmini0573-energyinjoules, def:physics:phyx-mini-0573:phyxminiproblems-problemphyxmini0573-speedinmeterspersecond, def:physics:phyx-mini-0573:phyxminiproblems-problemphyxmini0573-deuteriumfusionbombsetup}
  Positivity and nondegeneracy conditions for the physical branch
\end{definition}
\begin{proof}
  This is the declared model definition of Has Physical Fusion Parameters.
\end{proof}
\begin{definition}[reactant Rest Mass In Kilograms]
  \label{def:physics:phyx-mini-0573:phyxminiproblems-problemphyxmini0573-reactantrestmassinkilograms}
  \lean{PhyXMiniProblems.ProblemPhyXMini0573.reactantRestMassInKilograms}
  \uses{def:physics:phyx-mini-0573:phyxminiproblems-problemphyxmini0573-massinkilograms, def:physics:phyx-mini-0573:phyxminiproblems-problemphyxmini0573-nuclearspecies, def:physics:phyx-mini-0573:phyxminiproblems-problemphyxmini0573-nuclearreaction, def:physics:phyx-mini-0573:phyxminiproblems-problemphyxmini0573-deuteriumfusionbombsetup}
  Total rest mass on the reactant side of a channel, in kilograms
\end{definition}
\begin{proof}
  This is the declared model definition of reactant Rest Mass In Kilograms.
\end{proof}
\begin{definition}[product Rest Mass In Kilograms]
  \label{def:physics:phyx-mini-0573:phyxminiproblems-problemphyxmini0573-productrestmassinkilograms}
  \lean{PhyXMiniProblems.ProblemPhyXMini0573.productRestMassInKilograms}
  \uses{def:physics:phyx-mini-0573:phyxminiproblems-problemphyxmini0573-massinkilograms, def:physics:phyx-mini-0573:phyxminiproblems-problemphyxmini0573-nuclearspecies, def:physics:phyx-mini-0573:phyxminiproblems-problemphyxmini0573-nuclearreaction, def:physics:phyx-mini-0573:phyxminiproblems-problemphyxmini0573-deuteriumfusionbombsetup}
  Total rest mass on the product side of a channel, in kilograms
\end{definition}
\begin{proof}
  This is the declared model definition of product Rest Mass In Kilograms.
\end{proof}
\begin{definition}[Satisfies Fusion Yield Laws]
  \label{def:physics:phyx-mini-0573:phyxminiproblems-problemphyxmini0573-satisfiesfusionyieldlaws}
  \lean{PhyXMiniProblems.ProblemPhyXMini0573.SatisfiesFusionYieldLaws}
  \uses{def:physics:phyx-mini-0573:phyxminiproblems-problemphyxmini0573-massinkilograms, def:physics:phyx-mini-0573:phyxminiproblems-problemphyxmini0573-energyinjoules, def:physics:phyx-mini-0573:phyxminiproblems-problemphyxmini0573-speedinmeterspersecond, def:physics:phyx-mini-0573:phyxminiproblems-problemphyxmini0573-deuteriumfusionbombsetup, def:physics:phyx-mini-0573:phyxminiproblems-problemphyxmini0573-reactantrestmassinkilograms, def:physics:phyx-mini-0573:phyxminiproblems-problemphyxmini0573-productrestmassinkilograms}
  Governing yield relations for the independently stored quantities: reacted mass is the stated fraction of the initial fuel mass; reaction mass defect is reactant rest mass minus product rest mass; each event releases Δm c²; event count is reacted mass divided by rest mass consumed per event; total energy is event count times energy per event; and megaton rating is total energy divided by one megaton of TNT. No field assigns the requested rating or any displayed answer value
\end{definition}
\begin{proof}
  This is the declared model definition of Satisfies Fusion Yield Laws.
\end{proof}
\begin{definition}[Answer Choice]
  \label{def:physics:phyx-mini-0573:phyxminiproblems-problemphyxmini0573-answerchoice}
  \lean{PhyXMiniProblems.ProblemPhyXMini0573.AnswerChoice}
  Labels printed beside the four candidate fusion ratings
\end{definition}
\begin{proof}
  This is the declared model definition of Answer Choice.
\end{proof}
\begin{definition}[rating Megatons TNT]
  \label{def:physics:phyx-mini-0573:phyxminiproblems-problemphyxmini0573-answerchoice-ratingmegatonstnt}
  \lean{PhyXMiniProblems.ProblemPhyXMini0573.AnswerChoice.ratingMegatonsTNT}
  \uses{def:physics:phyx-mini-0573:phyxminiproblems-problemphyxmini0573-answerchoice}
  Megatons-of-TNT value printed beside each answer label
\end{definition}
\begin{proof}
  This is the declared model definition of rating Megatons TNT.
\end{proof}
\begin{definition}[recorded Answer Choice]
  \label{def:physics:phyx-mini-0573:phyxminiproblems-problemphyxmini0573-recordedanswerchoice}
  \lean{PhyXMiniProblems.ProblemPhyXMini0573.recordedAnswerChoice}
  \uses{def:physics:phyx-mini-0573:phyxminiproblems-problemphyxmini0573-answerchoice}
  Dataset metadata records answer C; this is not a theorem premise
\end{definition}
\begin{proof}
  This is the declared model definition of recorded Answer Choice.
\end{proof}
\begin{definition}[Matches Displayed Rating]
  \label{def:physics:phyx-mini-0573:phyxminiproblems-problemphyxmini0573-matchesdisplayedrating}
  \lean{PhyXMiniProblems.ProblemPhyXMini0573.MatchesDisplayedRating}
  \uses{def:physics:phyx-mini-0573:phyxminiproblems-problemphyxmini0573-answerchoice}
  Agreement with a rating displayed to two decimal places: the calculated value is within half a unit in the last displayed decimal place
\end{definition}
\begin{proof}
  This is the declared model definition of Matches Displayed Rating.
\end{proof}
% --- Archon declaration topology end ---
% --- Archon physics formalization source end ---
